\section{Conclusion}
\label{sec:conclusion}

This paper has shown that the choice of evaluation condition matters substantially for
creative work assessment---and that the most practically consequential finding is a warning,
not a capability claim. All quantitative results reported here are valid within the authors'
constructed evaluation framework (author-designed rubric, author-assigned tiers,
author-calibrated thresholds); they characterize discriminative power within that framework
and should not be read as validated against an external human expert standard.

Three findings stand up clearly within that framework.

First, the \emph{calibration inversion}: rubric-only evaluation (C1) inverts the expected
tier hierarchy, producing a pass rate for Standard Hunt puzzles that exceeds the pass rate
for MIT/Elite-tier puzzles (80\% versus 40\%). This is consistent across 15 real published
puzzles. The mechanism is traceable: elite puzzles deliberately sacrifice Clarity and
Solvability for Elegance and Reading Reward, and an equal-weight rubric treats that
trade-off as a defect. Expert identity, encoded in the full profile, reintroduces the
calibration that converts raw dimension scores into tier-appropriate verdicts.
\textbf{For practitioners: a bare rubric without expert identity context systematically
disadvantages your most ambitious work.}

Second, the \emph{feedback quality gradient}: C0 and C1 produce zero named evaluative
frameworks across all fifteen puzzles; C6 produces six---\emph{a-ha economy},
\emph{load-bearing test}, \emph{world-model consistency}, \emph{social fabric},
\emph{perceptual-shift}, \emph{visual grammar}---none of which appears in any input text
provided to the reviewers in those sessions.
(These counts were derived from two-pass manual coding of review transcripts by the research team; inter-rater reliability was not assessed.) Whether this reflects profile-driven emergence
or retrieval from the model's training prior is not established by this study. What is
established is that these frameworks appear in C6 reviews, not in C0/C1 reviews, and that
they identify specific, repair-pointing issues that rubric vocabulary does not surface.

Third, the \emph{Riven Standard as discriminator}: adding the Riven Standard as a scored
seventh dimension---with Elegance and Reading Reward double-weighted as proxies---achieves
a 24.5 percentage-point tier-separation gap between the MIT/Elite and Games Magazine tiers,
compared to 10.9 points for the equal-weight baseline. The Riven Standard~\cite{miller2016riven}
names whether a creative artifact's mechanism and its subject matter are structurally
inseparable; any creative domain can ask this question.

For practitioners building creative evaluation pipelines, the recommendation from this work
is compact. Use C6 profiles with the seven-dimension weighted rubric as the default
evaluation condition; it provides the most reliable verdicts and the richest actionable
feedback. Supplement with a C5 pass (lens only, no biography) for mechanical defect
detection, which catches visual and extraction errors that higher-level evaluation frameworks
do not surface. Add three non-redundant principles---\emph{Verify the Last Mile},
\emph{Blame the Player}, \emph{Snyder's Computer Test}---if the deepest available analysis
is required. Avoid C1 (rubric only) for any creative work with dimensional trade-offs by design; it
will invert your quality hierarchy and systematically disadvantage your most ambitious
artifacts.

Several questions are not resolved by this study and warrant dedicated future work:

\begin{itemize}

\item \emph{Human calibration.} The panel's verdicts have not been compared against
assessments by the named human practitioners on the same artifacts. Without a calibration
study on a shared puzzle set, the system's output can be characterized as internally coherent
and domain-grounded but not as validated against the human expert standard it simulates.
Closing this gap is the most direct path to practical deployment in high-stakes evaluation
contexts~\cite[cf.][]{dellaLibera2025calibration}.

\item \emph{Authoring symmetry.} This study examines whether reviewer profile quality
affects evaluation quality; an open question is whether authoring profile quality affects
creation quality by the same mechanism. If the same profile depth that makes a reviewer
more specific also makes a constructor more deliberate, the architecture has implications
for AI-assisted creative production, not just assessment.

\item \emph{Cross-domain generalization.} The evidence that rich profiles improve AI panel
performance in puzzle hunt design does not straightforwardly generalize to domains with
less codified quality criteria or less prominent expert communities. Level design,
interactive fiction, tabletop game mechanics, and musical composition are plausible
candidates, but it is an open question how much of the present methodology transfers
without modification.

\item \emph{Profile design as a research problem.} Which sections of a profile are
necessary for calibration, which are sufficient for framework emergence, and whether
profiles can be constructed to surface specific evaluative frameworks on demand are
questions that bear on both the efficiency and the reliability of the system.

\end{itemize}

The 29 domain-specific reviewer profiles (C8 format, $\sim$75 lines; approximately 1,800--2,200 tokens per profile, versus 2,800--3,400 tokens for the full C6 profiles), the seven-dimension
weighted rubric with anchor examples, all review transcripts from the eight-condition
ablation, and the analysis scripts used to compute tier separations and framework counts
are available at \url{[URL]}. We hope the profile architecture, the Riven Standard rubric,
and the calibration inversion result are useful to researchers working on AI-mediated
evaluation in other creative domains: the infrastructure is designed to be adapted, and the
failure modes documented here are not specific to puzzle hunts.
